% Copyright (c) 2026 Benjamin Stanley Frohman. Apache-2.0.
% Author: Frohman, Benjamin Stanley.
\documentclass[11pt]{article}
\usepackage[margin=1in]{geometry}
\usepackage{amsmath,amssymb,amsthm,lmodern}
\usepackage[T1]{fontenc}
\usepackage{xurl,hyperref,booktabs,tabularx}
\hypersetup{pdftitle={Degree-five Independence and a Nonzero Explicit Jacobian Class},pdfauthor={Frohman, Benjamin Stanley}}
\newcommand{\lean}[1]{{\upshape\path{#1}}}
\newtheorem{theorem}{Theorem}
\newtheorem{lemma}{Lemma}
\emergencystretch=2em
\title{\textbf{Degree-five Independence and a Nonzero Explicit Jacobian Class}}
\author{Frohman, Benjamin Stanley\\\small Independent Researcher, Texas, United States\\\small ORCID 0009-0006-7068-3718\quad\texttt{frohmanbenjamin@gmail.com}}
\date{Version 1.1, November 10, 2026 (v1.0: November 6, 2026)}
\begin{document}
\maketitle
\begin{abstract}
For a field $K$ and $k\ne0$, four specified degree-five monomials are proved
linearly independent in the actual quotient
$K[x,y]/(5kx^4y,kx^5+6y^5)$. The proof uses coefficient functionals annihilating
every element of the full ideal, with arbitrary polynomial multipliers. In
characteristic zero, the previously checked section identity and exact
coefficient certificates then imply that the actual quotient class of $HW$
is nonzero at every root of the specified modulus. This is an algebraic
milestone, not a geometric residue, Hodge or cycle-class theorem.
\end{abstract}

\section{The actual quotient and the claim}
Put $f_k=kx^5y+y^6$ and
\[
 I_k=(\partial_x f_k,\partial_y f_k)=(5kx^4y,kx^5+6y^5),\qquad
 Q_k=K[x,y]/I_k.
\]
Brackets denote the quotient map. The functions \lean{binarySextic},
\lean{jacobianIdeal}, \lean{Jacobian}, \lean{qx} and \lean{qy} in
\lean{Hodge/ExplicitJacobianQuotient.lean} are these actual objects, not an
abstract interface with a postulated basis.

\begin{theorem}\label{thm:independence}
For every field $K$ and $k\in K\setminus\{0\}$, the classes
\[
 [y^5],\quad[xy^4],\quad[x^2y^3],\quad[x^3y^2]
\]
are linearly independent over $K$ in $Q_k$.
\end{theorem}
No nontriviality of $Q_k$, no basis hypothesis and no geometric realization are
assumed. The theorem concerns four classes of degree five, not a basis of the
entire quotient. It does not assert a formalized spanning or dimension theorem
for a graded component either.

\section{Coefficient functionals on the full ideal}
For a polynomial $p$, let $c_{a,b}(p)$ be its coefficient of $x^ay^b$.
Define
\[
 L_0(p)=k c_{0,5}(p)-6c_{5,0}(p),\qquad
 L_i(p)=c_{i,5-i}(p)\quad(1\leq i\leq3).
\]
The minus sign and the $x^5$ coefficient are essential: the raw coefficient of
$y^5$ does not annihilate $kx^5+6y^5$ in characteristic zero.

\begin{lemma}\label{lem:ideal}
For every commutative coefficient ring, every $k$ and every $p\in I_k$, one has
$L_0(p)=L_1(p)=L_2(p)=L_3(p)=0$.
\end{lemma}
\begin{proof}
Write $p=A(5kx^4y)+B(kx^5+6y^5)$, where $A,B$ are arbitrary polynomials,
with no degree restriction. Coefficient multiplication gives
\begin{align*}
 c_{1,4}(p)&=c_{2,3}(p)=c_{3,2}(p)=0,\\
 c_{0,5}(p)&=6c_{0,0}(B),\qquad
 c_{5,0}(p)=k c_{0,0}(B).
\end{align*}
For example, the generator $x^4y$ cannot contribute at any of the first three
positions because its $x$ exponent is already too large. It cannot contribute
to either pure-power position. At the two pure powers, only the constant term
of $B$ contributes; every other multiplier monomial would increase an
exponent. Thus $L_0(p)=6k c_{0,0}(B)-6k c_{0,0}(B)=0$.
This proves the claim for the full ideal, not just its listed generators or
constant linear combinations thereof.
\end{proof}

In Lean, the proof uses \lean{Ideal.mem_span_pair} and
\lean{MvPolynomial.coeff_mul_monomial'}. These handle arbitrary ideal
multipliers exactly. The finite exponent checks do not postulate a homogeneous
basis or a Groebner basis.

\begin{proof}[Proof of Theorem~\ref{thm:independence}]
Suppose $[\Phi(c)]=0$, where
\[
 \Phi(c)=c_0y^5+c_1xy^4+c_2x^2y^3+c_3x^3y^2.
\]
Then $\Phi(c)\in I_k$. Lemma~\ref{lem:ideal} gives $c_1=c_2=c_3=0$ and
$k c_0=0$. Since $K$ is a field and $k\ne0$, also $c_0=0$.
This is precisely linear independence in the actual quotient.
\end{proof}

\section{Application to the actual explicit section}
For this application let $K$ be a characteristic-zero field equipped with its
specified $\mathbb Q$-algebra structure, and suppose
\begin{align*}
 P(k)={}&k^{24}+66081312000k^{18}-4811512860000k^{12}\\
        &+1167566400000k^6+583200000=0.
\end{align*}
The nonzero constant term implies $k\ne0$; this is a previously checked theorem,
not an extra nontriviality assumption.

Use the actual coefficient functions $b,c,d,e,g$ from
\lean{Hodge/ExplicitCurve.lean}, and set
\[
 H=x^2+bxy+cy^2,\qquad
 G=x^3+dx^2y+exy^2+gy^3,\qquad W=H_xG_y-H_yG_x.
\]
Their exact rational-polynomial definitions, together with every reduction and
Bezout certificate, are included in the pinned proof sources. In particular,
the computation is reproducible without numerical root approximation.
Writing $W=w_3x^3+w_2x^2y+w_1xy^2+w_0y^3$, the checked formulas are
\begin{align*}
 w_3&=2d-3b=k,& w_2&=4e-bd-6c,\\
 w_1&=6g+be-4cd,& w_0&=3bg-2ce.
\end{align*}
In the actual quotient, $[x^4y]=0$ and $k[x^5]+6[y^5]=0$. Expanding $HW$
and using these relations gives the already checked identity
\[
 [HW]=u_0[y^5]+u_1[xy^4]+u_2[x^2y^3]+u_3[x^3y^2],
\]
where
\begin{align*}
 u_0&=cw_0-6,& u_1&=bw_0+cw_1,\\
 u_2&=w_0+bw_1+cw_2,& u_3&=w_1+bw_2+ck.
\end{align*}
The $-6$ is the contribution of $k[x^5]$, not an omitted or freely chosen term.

The earlier exact polynomial certificates prove $u_0(k)\ne0$ whenever $P(k)=0$.
Their mechanism is $u_0=QP+r$ and $Ar+BP=1$ in $\mathbb Q[T]$, with the
explicit polynomials recorded in \lean{ExplicitResidueCertificates.lean}.
Evaluating at a root gives $A(k)u_0(k)=1$. This is an actual checked polynomial
identity, not a numerical test or a new coefficient-nonvanishing hypothesis.

\begin{theorem}
For every such field and every root $k$ of $P$, the actual class $[HW]\in Q_k$
is nonzero.
\end{theorem}
\begin{proof}
If $[HW]=0$, Theorem~\ref{thm:independence} and the displayed actual quotient
identity imply $u_0(k)=0$, contrary to the exact certificate. In fact only this
one coefficient certificate is needed here, although the earlier stage proves
all four nonzero.
\end{proof}

This statement is nonvacuous: the checked modulus has degree $24$ and has an
actual complex root. The new existential theorem supplies such a root together
with its nonzero $[HW]$, without assuming that the quotient is nontrivial.

\section{Formalization, reproduction and boundary}
\subsection*{Lean source and exact versions}
The module is \lean{Hodge/ExplicitJacobianIndependence.lean} in
\url{https://github.com/BenFrohman/HODGE}, code commit
\lean{52325b02f966139c99015443ca5fba8a725ef3bb} (pull request~\#49), on top of
\lean{main} at \lean{b13f57f}. Toolchain \lean{leanprover/lean4:v4.22.0};
Mathlib tag \lean{v4.22.0}, locked revision
\lean{79e94a093aff4a60fb1b1f92d9681e407124c2ca} (\lean{lake-manifest.json}).
The full source is in the repository; nothing here is a summary or certificate in
place of source. The module is imported by \lean{Hodge.lean}, so the default
library build checks it, and the CI sorry scan covers it.

\subsection*{Lean hypotheses versus the assumptions of this note}
\noindent\small
\begin{tabularx}{\linewidth}{@{}p{0.42\linewidth}X@{}}\toprule
Lean hypothesis & Assumption in this note\\\midrule
\texttt{[Field K]}, \texttt{(hk : k $\ne$ 0)} & Theorem~\ref{thm:independence}: $K$ a field, $k\ne0$.\\
\lean{jacobianIdeal k} (Stage~3 definition) & $I_k=(5kx^4y,\,kx^5+6y^5)$, the full ideal.\\
\texttt{[CharZero K] [Algebra $\mathbb{Q}$ K]} & ``characteristic zero''. For a field of characteristic~0 a $\mathbb Q$-algebra structure exists and is unique, so this adds nothing.\\
\texttt{hm : k\^{}24+66081312000*k\^{}18-\dots=0} & $P(k)=0$, the degree-24 modulus.\\
none (existential) & \lean{exists_complex_nonzero_HW_class}: a complex root of $P$ with $[HW]\ne0$.\\
\bottomrule\end{tabularx}\normalsize

\noindent No hypothesis is equivalent to its conclusion, and none is uninhabited:
$k=1$ over $\mathbb Q$ satisfies Theorem~\ref{thm:independence}, and the last row
exhibits a root for the second theorem.

\subsection*{One-command rebuild and axiom output}
From the repository root at the commit above:
\begin{verbatim}
lake exe cache get && lake build && python3 scripts/check_residue_axioms.py
\end{verbatim}
GitHub Actions ran exactly this (workflow \lean{certify.yml}, run
\url{https://github.com/BenFrohman/HODGE/actions/runs/42922528702}, passed;
library build and sorry check: run
\url{https://github.com/BenFrohman/HODGE/actions/runs/42922528652}, passed).
Its axiom output for this module, one line per theorem in source order:
\begin{verbatim}
'Hodge.ExplicitJacobianIndependence.monomial_exponent'
'Hodge.ExplicitJacobianIndependence.partial_x_monomial'
'Hodge.ExplicitJacobianIndependence.partial_y_monomial'
'Hodge.ExplicitJacobianIndependence.jacobian_ideal_coefficient_constraints'
'Hodge.ExplicitJacobianIndependence.representative_coefficient'
'Hodge.ExplicitJacobianIndependence.representative_x5_coefficient'
'Hodge.ExplicitJacobianIndependence.representative_mem_ideal_iff'
'Hodge.ExplicitJacobianIndependence.representative_class_formula'
'Hodge.ExplicitJacobianIndependence.representative_class_sum'
'Hodge.ExplicitJacobianIndependence.representative_class_eq_zero_iff'
'Hodge.ExplicitJacobianIndependence.degreeFiveClasses_linearIndependent'
'Hodge.ExplicitJacobianIndependence.section_HW_class_ne_zero'
'Hodge.ExplicitJacobianIndependence.exists_complex_nonzero_HW_class'
  each: depends on axioms: [propext, Classical.choice, Quot.sound]
PASS: 78 incremental theorem declarations; standard axioms only
\end{verbatim}
The root audit in the same run printed
\lean{audited 699 declarations, 699 with standard axioms only}.
No \lean{sorry}, \lean{admit}, \lean{native\_decide} or new \lean{axiom} occurs.
A kernel check shows that the stated Lean propositions follow from the stated
definitions; it does not by itself establish any geometric statement.

No geometric cycle-class map, residue-to-cohomology comparison, Hodge-type
identification, descent argument, or geometric counterexample is proved by this
milestone. In particular a nonzero class in this binary polynomial quotient
must not be relabeled a nonalgebraic Hodge class. The remaining geometric
bridges require their actual definitions and proofs.

\smallskip\noindent Copyright (c) 2026 Benjamin Stanley Frohman.
Existing Apache-2.0 licensing and third-party notices are preserved.
\end{document}
